% Notes of the Montreal seminar, 1970 (cote 129, typed part): archive cote 129, site page 25 (PDF page 26).
% typescript. Draft compared: none.
% AI-transcribed by Claude Opus 5.5 (claude-opus-5-5, adaptive thinking, effort high), Message Batches API, batch mtl129_full (msgbatch_01AGnDKHwy6ATCM7ehUa6263), 2026-11-13; not yet checked by eye.
\scanpage{129/25}
\diagram{\begin{tikzcd}
X \arrow[r, "\mathbb{F}_X"] \arrow[d, "g"'] & X \arrow[d, "g"'] \\
Y \arrow[r, "\mathbb{F}_Y"] & Y \quad .
\end{tikzcd}}

On obtient donc un endomorphisme $\mathbb{F}$ du foncteur identique de $\mathrm{Sch}/\mathbb{F}_p$ .

5.2. Soit $S$ un schéma fixé et $X$ un S-schéma. On note $X^{(p/S)}$
ou simplement $X^{(p)}$ le S-schéma image réciproque de $X$ par le change-%
ment de base $\mathbb{F}_S : S \to S$ .

Le diagramme ci-dessus étant commutatif, il existe une flè-%
che unique $\mathbb{F}_{X/S}$ telle que le diagramme suivant soit commutatif :

\diagram{\begin{tikzcd}
X \arrow[dr] \arrow[drr, "\mathbb{F}_X"] \arrow[ddr, "\mathbb{F}_{X/S}"'] & & \\
 & X^{(p)} \arrow[r] \arrow[d] & X \arrow[d] \\
 & S \arrow[r, "\mathbb{F}_S"] & S \quad .
\end{tikzcd}}

On appellera le morphisme $\mathbb{F}_{X/S}$ ainsi défini le \uline{morphisme de Frobenius}
\uline{relatif de} $X$ \uline{au-dessus de} $S$ .

L'application $X \to \mathbb{F}_{X/S}$ est un morphisme fonctoriel du fonc-%
teur identique de $\mathrm{Sch}/S$ dans le foncteur $X \to X^{(p/S)}$ . Ce dernier fonc-%
teur commute aux limites projectives finies, donc si $X$ est un S-groupe,
$X^{(p/S)}$ l'est aussi et, $\mathbb{F}_{X/S}$ étant fonctoriel, est un morphisme de groupes.

La construction de $X^{(p/S)}$ commute aux changements de base,
c'est-à-dire que pour tout morphisme $T \to S$ , on a

\mspage{\hand{25}}
